\section{Streaming Pipeline: Eventsim → Kafka → Flink → Snowflake}

This section walks through the three active stages of the real-time pipeline:
event generation, message transport, and stream processing.

\subsection{Event Simulation with Eventsim}

Eventsim is an open-source event generator that simulates the behaviour of users on a
fictional music streaming website.
It models realistic user sessions — login, browse, play tracks, skip songs, logout —
and emits each interaction as a JSON event to a Kafka topic.

\begin{itemize}
  \item Song metadata is drawn from a \textbf{10,000-song subset} of the Million Songs Dataset,
        providing realistic artist names, song titles, and durations.
  \item User demographics (age, location, subscription tier) are synthetically generated
        using US census data, giving the events geographic and behavioural diversity.
  \item Three topics are produced simultaneously:
    \begin{itemize}
      \item \texttt{listen\_events} — artist, song, duration, timestamp, user ID, location, subscription level.
      \item \texttt{page\_view\_events} — page name, HTTP method, status code, session ID, user context.
      \item \texttt{auth\_events} — session ID, success/failure flag, subscription level, timestamp.
    \end{itemize}
  \item The Confluent Schema Registry (port \texttt{8081}) enforces JSON schemas on each topic,
        preventing schema drift from corrupting downstream consumers.
\end{itemize}

\subsection{Flink Streaming Job}

The PyFlink job (\texttt{flink\_streaming/jobs/stream\_events.py}) consumes all three topics in
a single Flink execution graph using the Table API and Flink SQL.

\subsubsection{Schema and Source Registration}

Each topic's JSON payload is mapped to a set of Flink SQL column definitions in \texttt{schema.py}.
The \texttt{create\_kafka\_source()} function registers a temporary Kafka source table via DDL:

\begin{lstlisting}[language=SQL, caption={Generated Kafka source DDL for \texttt{listen\_events}.}]
CREATE TEMPORARY TABLE src_listen_events (
    `artist`        STRING,
    `song`          STRING,
    `duration`      DOUBLE,
    `ts`            BIGINT,
    `userId`        BIGINT,
    `level`         STRING,
    `city`          STRING,
    `state`         STRING,
    -- ...
) WITH (
    'connector'                    = 'kafka',
    'topic'                        = 'listen_events',
    'properties.bootstrap.servers' = 'broker:29092',
    'properties.group.id'          = 'flink-streamify',
    'scan.startup.mode'            = 'earliest-offset',
    'format'                       = 'json',
    'json.ignore-parse-errors'     = 'true'
)
\end{lstlisting}

\subsubsection{Transformation and Sink}

For each topic, the job enriches the raw stream and writes to a Snowflake JDBC sink:

\begin{itemize}
  \item \textbf{Timestamp conversion} — the raw \texttt{ts} field (epoch milliseconds) is converted
        to a proper \texttt{TIMESTAMP(3)} via \texttt{TO\_TIMESTAMP\_LTZ(ts / 1000, 3)}.
  \item \textbf{Partition columns} — \texttt{year\_}, \texttt{month\_}, \texttt{day\_}, \texttt{hour\_}
        are derived from the timestamp for efficient Snowflake partitioning.
  \item \textbf{String decoding UDF} — artist and song names in \texttt{listen\_events} and
        \texttt{page\_view\_events} are stored in a latin1 unicode-escape encoding.
        A Python UDF \texttt{string\_decode} converts them to UTF-8 before writing.
  \item \textbf{Processing timestamp} — \texttt{CURRENT\_TIMESTAMP} is appended as \texttt{processed\_time}
        to track pipeline latency.
\end{itemize}

All three INSERT statements are registered in a single \texttt{StatementSet} and submitted as
\textbf{one Flink job}, sharing one checkpoint cycle and one set of TaskManager slots:

\begin{lstlisting}[language=bash, caption={Submitting the PyFlink job after \texttt{docker compose up}.}]
docker exec -it flink-jobmanager \
    flink run -py /opt/flink/jobs/stream_events.py
\end{lstlisting}
